% Propuesta breve: tres avances principales de desarrollo
\subsubsection{Primera etapa: preparación de imágenes y desarrollo inicial}
En la primera etapa se desarrollaron herramientas para organizar imágenes, aplicar técnicas de aumento de datos y generar anotaciones mediante regiones delimitadoras. Posteriormente, la interfaz incorporó configuraciones por clase, persistencia de parámetros y un módulo inicial de clasificación de imágenes, estableciendo la base del flujo de trabajo visual.

\subsubsection{Segunda etapa: detección y entrenamiento}
La segunda etapa se centró en la identificación de múltiples objetos y el conteo a partir de imágenes. Se incorporaron procedimientos de entrenamiento de detectores, mecanismos de etiquetado y análisis de resultados, hasta adoptar modelos YOLO para el procesamiento de componentes de distintas clases.

\subsubsection{Tercera etapa: integración de Visbra con Kount-A}
En la última etapa se integraron los módulos de administración de perfiles y usuarios, las referencias ArUco para delimitar y medir la zona útil, la adquisición desde cámara y las secuencias de operación asociadas con el sistema embebido. Las revisiones posteriores incorporaron estimación de saturación visual, tiempos de espera configurables y mejoras en la presentación de resultados.

% Nota: los límites entre etapas constituyen una clasificación documental;
% verificar v4 y v7 si aparecen nuevos archivos originales.
